\section{Related Work}
\label{sec:related}
\noindent\textbf{Optimizing Quantum Circuit Simulation.} Previous studies on quantum circuit simulation utilize massive parallelism in GPUs and distributed clusters for full state vector simulation, distributing the state across aggregated VRAM in leadership-class supercomputers~\cite{10.1145/3771577, haner20175, chen201864}. Other studies propose static compilation and circuit partitioning to reduce communication overhead, maximizing kernel occupancy through precompiled kernels~\cite{xu2024atlas, zhang2021hyquas, li2021sv}. Additionally, tensor network-based approaches reduce memory footprint by decomposing quantum states into interconnected tensors, providing savings for low-entanglement circuits while facing exponential complexity in deep entangled circuits~\cite{burgholzer2022simulation, kim2025swiftn, huang2021efficient, lykov2022tensor}. Our study aligns with these efforts. However, existing server-centric approaches assume abundant memory and high-bandwidth interconnects, failing on edge devices with limited DRAM (e.g., 8~GB). In contrast, \EdgeQuantum's UMA-aware tiered memory architecture partitions the state vector across storage tiers via a UVM zero-copy data path. A 4-buffer ring pipeline and mantissa-truncated LZ4 ping-pong compression extend simulation capacity beyond DRAM.

\noindent\textbf{Memory Expansion and I/O Optimization Strategies.} Previous studies on memory expansion focus on extending simulation capacity by offloading state vectors to high-performance NVMe SSDs in server environments through OS-level demand paging or explicit memory mapping~\cite{park2022snuqs, zhang2025bmqsim}. Other work proposes lossless compression and adaptive encoding to trade CPU cycles for capacity~\cite{zhang2025bmqsim, wu2018memory, wu2019full, song2023efficient}. In the broader context of deep learning, prior studies reduce memory pressure by offloading intermediate feature maps to secondary storage~\cite{7783721, 264816, 10.1145/3630108}. While our study aligns with these techniques, generic strategies fail on edge devices due to SSD asymmetry and UMA constraints. Prior server-centric frameworks target high-end infrastructure and ignore variable-size update constraints on commodity edge storage. \EdgeQuantum's asynchronous I/O strategy serializes reads and writes to avoid NVMe contention, built on a UMA-aware UVM zero-copy data path that eliminates \texttt{cudaMemcpy} staging. A dual-file ping-pong mantissa-truncated LZ4 compression layer extends storage capacity beyond SSD limits.

\begin{comment}
  
\subsection{Optimizing Quantum Circuit Simulation}


There have been many studies that optimize quantum circuit simulation to enhance performance on high-end hardware.
Previous studies focused on utilizing massive parallelism in GPUs and distributed clusters for full state vector simulation.
These approaches accelerate execution and extend scalability by distributing the state vector across aggregated video memory (VRAM) in leadership-class supercomputers.
Other studies have proposed static compilation and circuit partitioning techniques to reduce communication overhead.
These methods analyze quantum circuits in advance to generate optimized execution plans or precompiled kernels, aiming to maximize kernel occupancy and minimize inter-node data transfer.
In addition, tensor network-based approaches have been proposed to reduce the memory footprint by decomposing quantum states into interconnected tensors.
These methods provide significant memory savings for shallow or low-entanglement circuits but face exponential complexity growth when simulating deep circuits with high entanglement.

Our study aligns with these prior efforts in improving the computational efficiency of quantum circuit simulation.
However, \EdgeQuantum aims to democratize quantum simulation by targeting resource-constrained edge devices rather than relying on inaccessible HPC infrastructure.
Existing datacenter-centric approaches such as \textit{ScaleQsim} and \textit{cuQuantum} assume abundant memory resources and high-bandwidth interconnects, which leads to execution failures on embedded devices where physical memory is strictly limited (e.g., 8 GB).
Through a tiered-memory architecture, \EdgeQuantum partitions the state vector across the storage hierarchy rather than across network nodes, and employs a specialized 4-buffer pipeline.
This allows \EdgeQuantum to enable large-scale simulation on commodity hardware without requiring millions of dollars in infrastructure.

\subsection{Memory Expansion and I/O Optimization Strategies}

To address the physical memory bottleneck, several simulation frameworks and system-level optimizations have been proposed to utilize secondary storage or compression.
Previous studies focused on extending simulation capacity by offloading state vectors to high-performance NVMe SSDs in server environments.
These approaches manage the movement of data pages between host memory and storage, often relying on OS-level demand paging or explicit memory mapping (mmap) to support simulations exceeding physical RAM capacity.
Other studies have proposed lossless compression and adaptive error-bounded encoding techniques to reduce memory consumption.
These methods dynamically compress state vectors during simulation, allowing systems to store larger quantum states within limited memory footprints by trading CPU cycles for capacity.
In the broader context of deep learning, approaches such as vDNN and FlashNeuron simplify memory management by offloading intermediate feature maps to backing stores.

These approaches highlight key techniques for improving capacity in data-intensive applications.
Similarly, \EdgeQuantum faces comparable challenges in edge quantum simulation, where the state vector size far exceeds the unified memory capacity of the SoC.
However, applying generic offloading techniques to edge devices is inefficient due to the bandwidth asymmetry of consumer-grade SSDs and the specific constraints of the Unified Memory Architecture (UMA).
To address this, \EdgeQuantum employs an asynchronous I/O strategy that serializes read and write operations to avoid NVMe contention, a critical optimization ignored by server-centric storage engines.
Combined with UVM-aware zero-copy buffering and adaptive mantissa-truncated LZ4 compression, \EdgeQuantum improves pipeline efficiency while minimizing CPU-GPU synchronization overhead.  
\end{comment}